% !TeX root = ../Lua-document-Template.tex 
% !TeX spellcheck = uk_UA

\chapter{Огляд літератури та теоретичні основи}

Тут ви пишете текст першого розділу. Абзацний відступ у 1.25 см сформується автоматично завдяки налаштуванням у головному файлі.

У цьому абзаці ми можемо послатися на авторитетне джерело інформації \cite{8509038}, яке автоматично потрапить до фінального списку літератури в кінці роботи. Поля сторінки (30-10-20-20 мм) та півтора інтервали будуть чітко дотримані на всіх сторінках, які компілюються.



\section{Перший розділ першої Глави}

Тест першого розділу Ось приклад цитування джерела: \cite{zhu2021deformable}.

		Це текст шрифтом Times New Roman. 

\textsf{Це текст шрифтом Arial.}

\texttt{Це моноширинний текст.}

\begin{figure}[htbp]
	\centering
	\includegraphics[width=0.8\textwidth]{images/Roof.jpg} % Ваше зображення
	\caption{Структурна схема алгоритму} % Підпис автоматично стане: Рисунок 1.1 — Структурна схема алгоритму
	\label{fig:algorithm} % Мітка для посилань у тексті
\end{figure}

Вставляння таблиці 
\begin{table}[htbp]
	\caption{Результати експериментальних досліджень} % Підпис має бути НАД таблицею
	\label{tab:results}
	\centering
	\begin{tabular}{|c|c|c|}
		\hline
		№ з/п & Параметр А & Параметр Б \\ \hline
		1     & 10.5       & 20.3       \\ \hline
		2     & 14.2       & 19.8       \\ \hline
	\end{tabular}
\end{table}
